% Section 6, printed pp. 52--61.
\section{The torsion component of the identity}\label{kps-sec-pictau}

\begin{definition}[Numerical and $\tau$-equivalence]\label{kps-def-numerical}
\dcref{6.1}\source{kleiman-picard:page-0052}
Line bundles $\LL,\NN$ are $\tau$-equivalent if some positive power of each
is algebraically equivalent. They are numerically equivalent if their
intersection numbers on every complete curve agree, equivalently if their
degrees agree on every integral curve or its normalization.
\end{definition}
\begin{definition}[Bounded family]\label{kps-def-bounded-family}
\dcref{6.2}\source{kleiman-picard:page-0052}
A family of invertible sheaves is bounded if it occurs among geometric fibers
of one invertible sheaf on $X\times T$ for a finite-type $k$-scheme $T$.
\end{definition}

\begin{theorem}[Characterizations of numerical triviality]\label{kps-thm-numerical}
\dcref{6.3}\source{kleiman-picard:page-0052}
Let $k$ be algebraically closed and $X$ projective. For an invertible $\LL$,
the following are equivalent: (a) $\LL$ is $\tau$-equivalent to $\OO_X$;
(b) $\LL$ is numerically equivalent to $\OO_X$; (c) the powers $\LL^p$ form a
bounded family; (d) $\chi(\mathcal F\ox\LL)=\chi(\mathcal F)$ for every coherent
$\mathcal F$; (e) this equality holds for every integral curve; and (f)
$\LL^p(1)$ is ample for every integer $p$. If $X$ is irreducible, these are
also equivalent to (g) $\chi(\LL^p(n))=\chi(\OO_X(n))$ for all $p,n$; if its
dimension is at least two, also to (h)
$\int c_1(\LL)h^{r-1}=\int c_1(\LL)^2h^{r-2}=0$.
\uses{kps-def-numerical,kps-def-bounded-family,kps-lem-section-bound,kps-lem-uniform-regularity,kps-prop-riemann,kps-prop-ample-positivity,kps-thm-hodge-index,kps-thm-snapper}
\end{theorem}
\begin{proof}
The implications are
$(c)\Rightarrow(a)\Rightarrow(d)\Rightarrow(e)\Longleftrightarrow(b)
\Rightarrow(c)\Rightarrow(f)\Rightarrow(b)$, together with
$(d)\Rightarrow(g)\Rightarrow(h)\Rightarrow(b)$ (and
$(g)\Rightarrow(b)$ for curves). If the powers of $\LL$ occur in a finite-type
family, two powers lie on one connected component, so their quotient is
algebraically equivalent to $\OO_X$; ampleness is open in a family, and a
common tensor power then gives (f). If (a) holds, replace $\LL$ by the power
which is algebraically equivalent to $\OO_X$. Constancy of Euler characteristic
in a connected family gives (d), hence (e) and (g). Riemann--Roch on every
integral curve gives $(e)\Longleftrightarrow(b)$, while (f) gives
\[
  0\le\deg(\LL^p(1)|_Y)=p\deg(\LL|_Y)+\deg(\OO_X(1)|_Y)
\]
for all integers $p$, forcing (b).

For irreducible $X$, expand the Hilbert polynomial of $\LL^p(n)$ in the
intersection numbers of $c_1(\LL)$ and $h=c_1(\OO_X(1))$. Equality (g)
annuls all terms involving $c_1(\LL)$, giving (h), and in dimension one it
directly gives degree zero. Conversely, for (h), reduce to the reduced case.
For dimension two use the Hodge index theorem. In higher dimension, choose a
general irreducible hyperplane section $H$ containing a prescribed integral
curve $Y$; the projection formula transfers the two vanishing intersections to
$H$, so induction gives numerical triviality on $H$ and hence on $Y$.

It remains to prove $(b)\Rightarrow(c)$. Lemma 6.6 supplies an $m$ for which
every numerically trivial $\NN$ is $m$-regular. Thus $\NN(m)$ is generated by
global sections, has dimension $M=\chi(\OO_X(m))$, and is a quotient of
$\OO_X(-m)^{\oplus M}$. The finite-type Quot scheme for this fixed Hilbert
polynomial carries a universal quotient. The open locus where that quotient is
invertible is still finite type (remove the proper image of the noninvertible
locus), and it contains every $\NN$; this is the required bounded family.
\end{proof}

\begin{lemma}[Uniform section bound]\label{kps-lem-section-bound}
\dcref{6.5}\source{kleiman-picard:page-0054}
For projective $r$-dimensional $X/k$ and coherent $\mathcal F$, there is
$B_{\mathcal F}$ such that for every numerically trivial $\LL$,
\[h^0(\LL\ox\mathcal F(n))\le B_{\mathcal F}\binom{n+r}{r}\quad(n\ge0).\]
\uses{kps-def-numerical,kps-prop-ample-positivity}
\end{lemma}
\begin{proof}
For $r=0$ take $B_{\mathcal F}=h^0(\mathcal F)$. Additivity in short exact
sequences reduces to $\mathcal F=\OO_X(p)$ on an integral $X$. If
$\LL(p)$ has a nonzero section, its effective divisor $D$ satisfies
\[
  0\le h^{r-1}[D]=h^{r-1}c_1(\LL)+p h^r=p h^r,
\]
so $p\ge0$ and $H^0(\LL(-1))=0$. For a hyperplane section $H$, the exact
sequence
\[
0\to\LL(n-1)\to\LL(n)\to\LL_H(n)\to0
\]
and induction on $r$ give
$h^0(\LL_H(n))\le B\binom{n+r-1}{r-1}$. Summing the successive
differences from the vanishing at $n=-1$ yields
$h^0(\LL(n))\le B\binom{n+r}{r}$. For $p<0$, $\OO_X(p)\subset\OO_X$;
for $p\ge0$, the elementary binomial product bound gives the same estimate
with $B_{\OO_X(p)}=B\binom{p+r}{r}$.
\end{proof}

\begin{lemma}[Uniform regularity and Hilbert polynomial]\label{kps-lem-uniform-regularity}
\dcref{6.6}\source{kleiman-picard:page-0055}
There is an integer $m$ such that every numerically trivial invertible sheaf on
projective $X/k$ is $m$-regular and
$\chi(\LL(n))=\chi(\OO_X(n))$ for every $n$.
\uses{kps-lem-section-bound,kps-def-numerical}
\end{lemma}
\begin{proof}
Induct on $r=\dim X$. The assertion is clear for $r=0$. Choose $q>0$ so
$\LL(q)$ is very ample, and choose effective divisors $F,G$ with
$\OO_X(F)=\OO_X(q)$ and $\OO_X(G)=\LL(q)$. Tensoring the two divisor exact
sequences by $\LL^p(n+q)$ gives
\[
0\to\LL^p(n)\to\LL^p(n+q)\to\LL_F^p(n+q)\to0,
\quad
0\to\LL^{p-1}(n)\to\LL^p(n+q)\to\LL_G^p(n+q)\to0.
\]
Induction on the hyperplane dimension shows
$\chi(\LL^p(n))=\phi_1(n)p+\phi_0(n)$. The section bound of Lemma 6.5,
together with Serre vanishing in degrees at least two, rules out a nonzero
$\phi_1$ (otherwise $h^0(\LL^p(n))$ is unbounded as $p\to\pm\infty$).
Thus $\chi(\LL^p(n))=\chi(\OO_X(n))$ for every $p,n$. Finally, Mumford's
regularity induction applies to the same exact sequences: the uniform section
bound replaces the ideal-sheaf estimate and yields one integer $m$ working for
all numerically trivial $\LL$.
\end{proof}

\begin{definition}[Torsion component]\label{kps-def-torsion-component}
\dcref{6.8}\source{kleiman-picard:page-0056}
For a group scheme $G/S$, let $\varphi_n$ be the $n$th power map and set
\[G^\tau=\bigcup_{n>0}\varphi_n^{-1}(G^0).\]
\end{definition}
\begin{lemma}[Properties of $G^\tau$]\label{kps-lem-torsion-component}
\dcref{6.9}\source{kleiman-picard:page-0056}
For a commutative locally finite-type group scheme over a field, $G^\tau$ is
an open subgroup and commutes with field extension. If quasi-compact, it is
closed and finite type.
\uses{kps-def-torsion-component,kps-lem-group-scheme}
\end{lemma}
\begin{proof}
Each inverse image $\varphi_n^{-1}(G^0)$ is open and closed and the family is
filtered; quasi-compactness reduces the union to finitely many terms.
\end{proof}

\begin{proposition}[The torsion component over a field]\label{kps-prop-pict-field}
\dcref{6.12}\source{kleiman-picard:page-0057}
If $\IPic_{X/k}$ exists and represents the fppf functor, then $\Pict_{X/k}$
is open and commutes with field extension; if $X$ is projective it is closed
and finite type.
\uses{kps-lem-torsion-component,kps-thm-numerical,kps-prop-lft}
\end{proposition}
\begin{proof}
The lemma reduces the assertion to quasi-compactness. After algebraic closure,
the final bounded-family construction in Theorem 6.3 covers all numerically
trivial sheaves, hence all of $\Pict$ by the equivalence with $\tau$-equivalence.
\end{proof}

\begin{remark}[Complete varieties and Jacobians]\label{kps-rmk-pict-field}
\dcref{6.14}\source{kleiman-picard:page-0057}
For complete, not necessarily projective, varieties the same finiteness follows
by Chow's lemma and pullback along a projective modification. For an integral
proper curve, numerical triviality is degree zero, so $\Picz=\Pict$.
\end{remark}

\begin{theorem}[Finiteness of the relative torsion component]\label{kps-thm-pict-finite}
\dcref{6.16}\source{kleiman-picard:page-0058}
If $X/S$ is flat, projective Zariski locally, with integral geometric fibers,
then $\Pict_{X/S}$ is an open and closed finite-type group scheme and commutes
with base change. If $X/S$ is projective and $S$ Noetherian, it is
quasi-projective.
\uses{kps-thm-main,kps-lem-torsion-component,kps-prop-pict-field}
\end{theorem}
\begin{proof}
Work locally where $X/S$ is projective. Fiberwise the torsion components are
open, commute with residue-field extension, and their union is the relative
torsion locus. On an affine chart, the numerical characterizations make its
preimage in an etale line-bundle family open and closed. The bounded-family
criterion and the uniform regularity/Quot construction prove quasi-compactness;
local finite type then gives finite type, and the lemma gives closedness.
\end{proof}

\begin{corollary}[Bounded component group]\label{kps-cor-component-group}
\dcref{6.17}\source{kleiman-picard:page-0059}
For Noetherian $S$ as above, the quotient
$\Pict_{X_{\bar k_s}/\bar k_s}(\bar k_s)/\Picz_{X_{\bar k_s}/\bar k_s}(\bar k_s)$ is
finite of uniformly bounded order.
\uses{kps-prop-pic0,kps-prop-pict-field,kps-thm-pict-finite}
\end{corollary}
\begin{proof}
It counts connected components of the finite-type torsion scheme. Base change
identifies the geometric fibers, and Noetherian induction bounds their number.
\end{proof}

\begin{remark}[Proper-family torsion and Neron--Severi bounds]\label{kps-rmk-pict-finite-general}
\dcref{6.19}\source{kleiman-picard:page-0059}
For a proper family whose Picard scheme represents the fppf functor, the
torsion component is open of finite type; over a Noetherian base the geometric
torsion quotient is finite of uniformly bounded order even when the total
Picard scheme is not known to exist. The corresponding Neron--Severi rank is
also finite and bounded, although its value need not be constructible in the
base point.
\end{remark}

\begin{theorem}[Fixed Hilbert polynomial loci]\label{kps-thm-fixed-polynomial}
\dcref{6.20}\source{kleiman-picard:page-0060}
For projective flat $X/S$ with integral geometric fibers and a polynomial
$\phi\in\QQ[n]$, the locus $\IPic^\phi_{X/S}$ of line bundles with Hilbert
polynomial $\phi$ is open and closed, finite type, disjointly covers the Picard
scheme, and commutes with base change. Over Noetherian $S$ it is
quasi-projective.
\uses{kps-thm-main,kps-prop-lft}
\end{theorem}
\begin{proof}
The Hilbert polynomial is locally constant in flat families, so the loci are
open and closed and cover. The same boundedness argument as for $\Pict$ is
simpler because the polynomial is fixed; finite type and quasi-projectivity
follow from the Quot/Hilbert parameter spaces.
\end{proof}

\begin{remark}[Numerical bounds for finite-type loci]\label{kps-rmk-fixed-polynomial-bounds}
\dcref{6.24}\source{kleiman-picard:page-0060}
For Noetherian $S$ and projective flat $X/S$ with integral $r$-dimensional
geometric fibers, a subset of the Picard scheme is finite type when the
coefficients $a_{r-1}$ of the fiber Hilbert polynomials are bounded above and
below and $a_{r-2}$ is bounded below. Equivalently, it suffices to bound
$\int\ell h^{r-1}$ above and below and $\int\ell^2h^{r-2}$ below. These
criteria follow by reducing to normal fibers and applying the fixed-polynomial
theorem; they are equivalent to finite-type representability for restriction
maps along suitable relative divisors.
\end{remark}

\begin{remark}[Curves and torsion-free degree components]\label{kps-rmk-curves}
\dcref{6.22}\source{kleiman-picard:page-0060}
For a family of integral curves, degree-$m$ pieces are open and closed finite
type, are fppf torsors under $\Picz$, and there is no torsion:
$\Picz_{X/S}=\Pict_{X/S}$. The same equality holds for Abelian schemes.
\end{remark}
\begin{corollary}[Connected components]\label{kps-cor-connected-components}
\dcref{6.25}\source{kleiman-picard:page-0061}
Connected components of $\IPic_{X/S}$ are open and closed finite-type schemes.
\uses{kps-thm-fixed-polynomial}
\end{corollary}
\begin{remark}[Need for global projectivity]\label{kps-rmk-connected-projectivity}
\dcref{6.26}\source{kleiman-picard:page-0061}
The projectivity hypothesis in the connected-components corollary cannot be
weakened to properness or projectivity only Zariski locally: a base curve with
two components meeting in two points can have a family projective near each
component whose Picard scheme has a connected component that is not of finite
type.
\end{remark}
\begin{corollary}[Power maps]\label{kps-cor-power-map}
\dcref{6.27}\source{kleiman-picard:page-0061}
For $n\ne0$, the power map $\varphi_n:\IPic_{X/S}\to\IPic_{X/S}$ is of finite
type.
\uses{kps-thm-fixed-polynomial,kps-cor-connected-components}
\end{corollary}
\begin{remark}[Proper-family power maps]\label{kps-rmk-power-map-general}
\dcref{6.28}\source{kleiman-picard:page-0061}
For a proper family with an fppf-represented Picard scheme, every nonzero power
map on the Picard scheme is of finite type; the projective case reduces to this
statement by the same descent and reduction steps used for the torsion theorem.
\end{remark}
